\documentclass[../main.tex]{subfiles}

\begin{document}
In the main chapters, we focused on cases where the oracles and the master-problem would be tailored to each game. While the ability to leverage the problem structure by providing custom oracles based on closed‑form best responses or reductions to classic combinatorial problems may be powerful, it naturally puts into question the applicability of our method.

It is certainly not possible to develop an oracle generic enough to handle all games, as this reduces to global optimization over arbitrary functions. Nevertheless, we will show that an off-the-shelf global optimization solver is sufficient to solve a surprisingly wide variety of continuous games automatically. To this end, we use the suite of continuous games compiled from the game‑theoretic and optimisation literature presented in Appendix \ref{appendix:games}. The collection spans multi-player polynomial, rational‑polynomial, and more generic games, with action spaces that include n-spheres and simplexes.

A consistent implementation also enables us to report meaningful and comparable runtimes across games. Our experiments show that a generic algorithm can solve many continuous games efficiently. Notably, our implementation locates the pure‑strategy Nash equilibrium of Example \ref{ex.zhou19.5.1} within a distance of less than $3\times 10^{-4}$ in 0.086 seconds compared to the 184.5 seconds reported originally.

\section{Experimental Evaluation}
We show the performance of \ref{alg:do} implemented in \emph{Julia 1.11} using \emph{Gurobi 12} as the oracle, and \emph{gambit-logit} as the master-problem solver. Experiments were performed on a desktop PC with \emph{AMD Ryzen 5 4600G} running \emph{Linux}. The reported solve times include the complete run of the algorithm after a warm-start to force compilation. The stopping criterion was uniformly set to $\epsilon=0.001$, i.e., no player should be able to deviate by more than $\epsilon$ utility units.

\subsection{Experiments on General Blotto}
The utility functions in \emph{General Blotto} games are sums of non-linearly weighted margins of victory over sets of fronts. The family of weighing functions described by \textcite{Golman2009} is \(f(z) = \operatorname{sgn}(z)\,|z|^p\), parametrized by $p$ which is strictly positive for continuous games. The margins of victory are differences of products of the allocations to each front.

While products, exponents, sums and differences should be available through most modeling interfaces for non-linear problems; neither $\operatorname{sgn}$, nor the absolute value is available via the non-linear input for \emph{Gurobi}. Instead, we modelled the absolute value as $\sqrt{z^2}$.

\begin{figure}
    \centering
    \begin{tikzpicture}
    \begin{axis}[
        xmin=-1, xmax=1,
        ymin=-1.2, ymax=1.2,
        xlabel=Margin of victory,
        legend pos=outer north east,
        colormap/Set1,
        legend cell align={left},
        width=4cm,
        height=4cm
    ]
    \addplot[index of colormap={3 of Set1}, samples=101]{sign(x)*abs(x)^2};
    \addlegendentry{$p=2.0$}

    \addplot[index of colormap={2 of Set1}]{x};
    \addlegendentry{$p=1.0$}

    \addplot[index of colormap={1 of Set1}, samples=301]{sign(x)*sqrt(abs(x))};
    \addlegendentry{$p=0.5$}

    \addplot[index of colormap={4 of Set1}, domain=-1:-0.0001] {-1};
    \addplot[index of colormap={4 of Set1}, domain=0.0001:1] {1};
    \addplot[index of colormap={4 of Set1}, only marks,mark=o] coordinates {(0,1)};
    \addplot[index of colormap={4 of Set1}, only marks,mark=o] coordinates {(0,-1)};
    \addplot[index of colormap={4 of Set1}, only marks,mark=*] coordinates {(0,0)};
    \addlegendentry{$p=0.0$}
    \end{axis}
    \end{tikzpicture}
    \caption{Behavior of weighing functions for \emph{General Blotto} for $p\in\{0, 0.5, 1, 2\}$.}
    \label{fig:blotto.weighing.functions}
\end{figure}

For each Example \ref{ex.golman09.isolated}--\ref{ex.golman09.sets}, we computed equilibria in a game with 5 fronts for each of the values $p\in\{0.5, 1, 2\}$ as representative cases of each behavior as shown in \cref{fig:blotto.weighing.functions}. The weighing functions for simplify to: $z\,\sqrt{z^2}$, for $p=2$; and $z$ for $p=1$.

In the smooth case of $p=2$, \emph{Double Oracle} was able to recover a pure equilibrium for the game with isolated fronts, and mixed equilibria in the other two games where pure equilibria do not exist. The finding an initial feasible solution approximately 2 milliseconds for each example. Initializing with a best response solution did not consistently improve the runtimes.
\begin{description}
\item[Result of Example \ref{ex.golman09.isolated}] In the case of 5 isolated fronts with $p=2$, we computed the pure equilibrium $x_1^\star = (1.0, 0.0, 0.0, 0.0, 0.0)$, $x_2^\star = (1.0, 0.0, 0.0, 0.0, 0.0)$
in a single iteration, with a total runtime of $0.4\,s$.
\item[Result of Example \ref{ex.golman09.pairs}] For the case of all pairs of front we computed the mixed $\epsilon$-NE:
\begin{align*}
\mu_1^\star &\approx \frac{1}{3}\delta_{ (1, 0, 0, 0, 0) } + \frac{1}{3}\delta_{ (0, 0, 0, 1, 0) } + \frac{1}{3}\delta_{ (0, 1, 0, 0, 0) },\\
\mu_2^\star &\approx \frac{1}{3}\delta_{ (1, 0, 0, 0, 0) } + \frac{1}{3}\delta_{ (0, 1, 0, 0, 0) } + \frac{1}{3}\delta_{ (0, 0, 0, 1, 0) }
\end{align*}
in 5 iterations and a time of $0.91\,s$.
Equilibrium:
\item[Result of Example \ref{ex.golman09.sets}] With all sets equally valued, we got the following mixed $\epsilon$-NE:
\begin{align*}
\mu_1^\star &\approx \frac{1}{3}\delta_{ (0, 0, 0, 0, 1) } + \frac{1}{3}\delta_{ (1, 0, 0, 0, 0) } + \frac{1}{3}\delta_{ (0, 1, 0, 0, 0) },\\
\mu_2^\star &\approx \frac{1}{3}\delta_{ (1, 0, 0, 0, 0) } + \frac{1}{3}\delta_{ (0, 0, 0, 0, 1) } + \frac{1}{3}\delta_{ (0, 1, 0, 0, 0) };
\end{align*}
in 5 iterations and a time of $2.2\,s$.
\end{description}

\emph{Double Oracle} was also able to successfully converge to pure equilibria in all the \emph{General Blotto} games with $p=1$:
\begin{description}
\item[Result of Example \ref{ex.golman09.isolated}] For isolated fronts, DO computed $\epsilon$-NE:
$x_1^\star \approx (1, 0, 0, 0, 0)$, $x_2^\star \approx (1, 0, 0, 0, 0)$
in a single iteration and a time of $0.004\,s$.
\item[Result of Example \ref{ex.golman09.pairs}] For all pairs, DO computed the following $\epsilon$-NE in two iterations and a time of  $0.018\,s$:
$x_1^\star \approx \left(\frac{1}{5}, \frac{1}{5}, \frac{1}{5}, \frac{1}{5}, \frac{1}{5}\right)$, $x_2^\star \approx \left(\frac{1}{5}, \frac{1}{5}, \frac{1}{5}, \frac{1}{5}, \frac{1}{5}\right)$
\item[Result of Example \ref{ex.golman09.sets}] For all sets, DO computed $\epsilon$:
$x_1^\star \approx \left(\frac{1}{5}, \frac{1}{5}, \frac{1}{5}, \frac{1}{5}, \frac{1}{5}\right)$, $x_2^\star \approx \left(\frac{1}{5}, \frac{1}{5}, \frac{1}{5}, \frac{1}{5}, \frac{1}{5}\right)$
in two iterations and $1.5\,s$.
\end{description}

The utility functions are continuous but not differentiable for $0<p<1$; Since \emph{Gurobi} does not accept neither the sign function nor absolute value through its non-linear interface. Attempting to model the utility as $\sfrac{z}{\sqrt[4]{z^2}}$ introduces a discontinuity which \emph{Gurobi} is able to avoid; however, the oracle did not find a best response within two minutes. The local solver \emph{Ipopt} failed to compute a solution as well, since the utility is not differentiable. The global solver \emph{Couenne} via \emph{AMPL} depends on \emph{Ipopt} and was similarly unsuccessful.

While our generic algorithm failed on all non-differentiable instances, specifying a custom oracle makes it possible to send more information about the utility to solvers. We investigated this approach in section \ref{sec:custom.oracles}.

To illustrate the complexity of the utility functions in \cref{ex.golman09.sets} where all sets of five fronts are equally important, we write it out partially below:
{
\begin{align*}
u_1(x_1,x_2)
    &= \operatorname{sgn}(x_{11}x_{13}x_{14} - x_{21}x_{23}x_{24})\abs{x_{11}x_{13}x_{14} - x_{21}x_{23}x_{24}}^p\\
    &+ \operatorname{sgn}(x_{11}x_{12}x_{15} - x_{21}x_{22}x_{25})\abs{x_{11}x_{12}x_{15} - x_{21}x_{22}x_{25}}^p\\
    &+ \operatorname{sgn}(x_{13}x_{14}x_{15} - x_{23}x_{24}x_{25})\abs{x_{13}x_{14}x_{15} - x_{23}x_{24}x_{25}}^p\\
    &+ \operatorname{sgn}(x_{11}x_{12}x_{13}x_{15} - x_{21}x_{22}x_{23}x_{25})\abs{x_{11}x_{12}x_{13}x_{15} - x_{21}x_{22}x_{23}x_{25}}^p\\
    &+ \operatorname{sgn}(x_{11}x_{13}x_{14}x_{15} - x_{21}x_{23}x_{24}x_{25})\abs{x_{11}x_{13}x_{14}x_{15} - x_{21}x_{23}x_{24}x_{25}}^p \\
    &+ \dots \text{26 more terms}.
\end{align*}
}
The Double Oracle algorithm still successfully converges for a game with $p=2$ and $40$ isolated fronts in $11$ seconds; with $20$ fronts when all pairs are important in $12.7$ seconds; and for a game with $7$ fronts, of which all sets are important in $23.6$ seconds, at which point the utility function contains the weighing function $127$ times.

\subsection{Experiments on Games by Parillo, Stein and Ozdaglar}
\begin{description}
\item[Result of Example \ref{ex.parrilo06.2.1}] DO converges to the $\epsilon$-NE:
\begin{align*}
\mu_1^\star &\approx 0.5\delta_{ -1.0 } + 0.5\delta_{ 1.0 },\\
\mu_2^\star &\approx 1.0\delta_{ 0.0 }.
\end{align*}
in $4$ iterations, and a time of $0.016\,s$.
\item[Result of Example \ref{ex.parrilo06.3.1}] After $0.036$ seconds and 6 iterations we get the pure equilibrium:
\[x_1^\star \approx 0.39, \; x_2^\star \approx 0.63.\]
\item[Result of Example \ref{ex.parrilo06.3.2}] After $0.073\,s$, and 9 iterations, DO converges to the $\epsilon$-NE:
\begin{align*}
\mu_1^\star &\approx 0.118\delta_{ 0.18 } + 0.882\delta_{ 0.2 },\\
\mu_2^\star &\approx 0.224\delta_{ -1.0 } + 0.776\delta_{ 1.0 }.
\end{align*}
\item[Result of Example \ref{ex.stein07.4.3.1}] After $0.018\,s$ and 4 iterations, DO converges to
$x_1^\star \approx 1$, $x_2^\star \approx 1$.
\item[Result of Example \ref{ex.stein08.2.3}]After $0.13\,s$, and 8 iterations, DO converges to the $\epsilon$-NE:
\begin{align*}
\mu_1^\star &\approx 0.056\delta_{ 0.111 } + 0.379\delta_{ 0.115 } + 0.565\delta_{ -1.0 },\\
\mu_2^\star &\approx 0.309\delta_{ 0.73 } + 0.691\delta_{ 0.71 }.
\end{align*}
\item[Result of Example \ref{ex.stein08.2.4}] It takes 22 iterations and $0.48\,s$ to converge to the $\epsilon$-NE:
\begin{align*}
\mu_1^\star &\approx 0.066\delta_{ -2.5 } + 0.069\delta_{ -3.0 } + 0.069\delta_{ -2.0 } + 0.078\delta_{ 2.78 } + 0.078\delta_{ -1.5 } \\ &+ 0.091\delta_{ 2.28 } + 0.091\delta_{ -1.0 } + 0.107\delta_{ 1.78 } + 0.107\delta_{ -0.5 } + 0.122\delta_{ 1.28 } + 0.122\delta_{ 0.0 },\\
\mu_2^\star &\approx 0.076\delta_{ -3.0 } + 0.076\delta_{ -2.5 } + 0.078\delta_{ 2.78 } + 0.078\delta_{ -2.0 } + 0.081\delta_{ 2.28 } + 0.081\delta_{ -1.5 } \\ &+ 0.085\delta_{ 1.78 } + 0.085\delta_{ -1.0 } + 0.089\delta_{ 1.28 } + 0.089\delta_{ -0.5 } + 0.092\delta_{ 0.78 } + 0.092\delta_{ 0.0 }.
\end{align*}
The game has much smaller known equilibria.
\item[Result of Example \ref{ex.stein08.3.10}] Finally, after 3 iterations and $0.017\,s$ we get the pure equilibrium:
\begin{align*}
x_1^\star \approx 1.0, \;
x_2^\star \approx 1.0, \;
x_3^\star \approx -0.5 .
\end{align*}
\end{description}


\subsection{Experiments on Games with Interesting Weaker Solution Concepts}
\begin{description}
\item[Result of Example \ref{ex.daskalakis.d}] DO converges in $0.015\,s$ and 1 iteration to the pure equilibrium:
\begin{align*}
x_1^\star \approx 0.5,\;
x_2^\star \approx 0.5.
\end{align*}
\item[Result of Example \ref{ex.daskalakis.e.1st}] DO converges in $0.063\,s$ and 4 iterations to the $\epsilon$-NE:
\begin{align*}
\mu_1^\star &\approx 0.5\delta_{ 1 } + 0.5\delta_{ -1 },\\
\mu_2^\star &\approx \delta_{ 0 }.
\end{align*}
\item[Result of Example \ref{ex.mertikopoulos18.fig1}]  DO converges in $0.087\,s$ and 6 iterations to the $\epsilon$-NE:
\begin{align*}
\mu_1^\star &\approx 0.435\delta_{ 0.44 } + 0.565\delta_{ 0.33 },\\
\mu_2^\star &\approx 0.192\delta_{ 2.5\times 10^{-6} } + 0.238\delta_{ 0 } + 0.569\delta_{ 1 }.
\end{align*}
\item[Result of Example \ref{ex.mertikopoulos18.2.2}] After 3 iterations and $0.027\,s$, DO converges to the equilibrium
\begin{align*}
x_1^\star \approx 0, \;
x_2^\star \approx 0.
\end{align*}
\item[Result of Example \ref{ex.razaviyayn20.5.1}] After 3 iterations and $0.016\,s$, DO converges to the $\epsilon$-NE:
\begin{align*}
\mu_1^\star &\approx 0.5\delta_{ -1 } + 0.5\delta_{ 1 },\\
\mu_2^\star &\approx 0.5\delta_{ -6.28 } + 0.5\delta_{ 6.28 }.
\end{align*}
\item[Result of Example \ref{ex.razaviyayn20.5.3}] After 5 iterations and $0.06\,s$, DO converges to the $\epsilon$-NE:
\begin{align*}
\mu_1^\star &\approx 0.333\delta_{ 1.0 } + 0.667\delta_{ -0.5 },\\
\mu_2^\star &\approx 0.312\delta_{ 1.0 } + 0.688\delta_{ -1.0 }.
\end{align*}
\item[Result of Example \ref{ex.zheng23.convex.nonconcave}]
After 9 iterations and $0.17\,s$, DO converges to the $\epsilon$-NE:
\begin{align*}
\mu_1^\star &\approx 0.064\delta_{ -0.04 } + 0.936\delta_{ -0.02 },\\
\mu_2^\star &\approx 0.03\delta_{ -0.05 } + 0.064\delta_{ 1 } + 0.302\delta_{ -0.03804 } + 0.302\delta_{ -0.03802 } + 0.302\delta_{ -0.03801 }.
\end{align*}
\item[Result of Example \ref{ex.zheng23.kl.nonconcave}]
After 2 iterations and $0.014\,s$, DO converges to the pure equilibrium:
\begin{align*}
x_1^\star \approx 0, \;
x_2^\star \approx 0.
\end{align*}
\item[Result of Example \ref{ex.zheng23.bilinearly.coupled.minimax}]
After 8 iterations and $0.18\,s$, DO converges to the $\epsilon$-NE:
\begin{align*}
\mu_1^\star &\approx 0.001\delta_{ 2.26 } + 0.024\delta_{ -2.25 } + 0.476\delta_{ -2.24 } + 0.499\delta_{ 2.24 },\\
\mu_2^\star &\approx 0.25\delta_{ 2.24 } + 0.25\delta_{ 2.24 } + 0.5\delta_{ -2.24 }.
\end{align*}
\item[Result of Example \ref{ex.zheng23.forsaken}]
After 5 iterations and $0.13\,s$, DO converges to the $\epsilon$-NE:
\begin{align*}
\mu_1^\star &\approx 0.498\delta_{ 1.32 } + 0.502\delta_{ -1.31 },\\
\mu_2^\star &\approx 0.33\delta_{ -1.31 } + 0.67\delta_{ 1.31 }.
\end{align*}
\end{description}

\subsection{Experiments on Two-Player Zero-Sum Polynomial Games}
The following examples are two-player zero-sum polynomial games, that originally served as experiments for the algorithm developed by \textcite{nie.yang.ea:saddle}. Their algorithm can prove the existence of (global) saddle points, i.e. pure Nash Equilibria.

Double Oracle usually converged to mixed strategies, even in the examples where pure equilibria exist. Tightening the $\epsilon$ did not result in simpler soultions.

\begin{description}
\item[Result of Example \ref{ex.nie21.6.1.i}]
After $0.077\,s$ and $7$ iterations, DO converged to the $\epsilon$-NE:
\begin{align*}
\mu_1^\star &\approx 0.001\delta_{ (0, 0, 1) } + 0.999\delta_{ (0, 1, 0) },\\
\mu_2^\star &\approx 0.002\delta_{ (0.54, 0.46, 0) } + 0.998\delta_{ (0, 0.49, 0.51) }.
\end{align*}
\item[Result of Example \ref{ex.nie21.6.1.ii}]
After $0.019\,s$ and $2$ iterations, DO converged to the $\epsilon$-NE:
\begin{align*}
x_{1}^\star \approx  (0, 0, 1) , \;
x_{2}^\star \approx  (0, 0, 1) .
\end{align*}
\item[Result of Example \ref{ex.nie21.6.1.iii}]
After $1.1\,s$ and $22$ iterations, DO converged to the $\epsilon$-NE:
\begin{align*}
\mu_1^\star &\approx 0.003\delta_{ (0.23, 0.24, 0.27, 0.26) } + 0.03\delta_{ (0.24, 0.24, 0.25, 0.27) } + 0.966\delta_{ (0.25, 0.25, 0.25, 0.25) },\\
\mu_2^\star &\approx 0.236\delta_{ (0, 0, 1, 0) } + 0.236\delta_{ (0, 0, 0, 1) } + 0.242\delta_{ (1, 0, 0, 0) } + 0.286\delta_{ (0, 1, 0, 0) }.
\end{align*}
\item[Result of Example \ref{ex.nie21.6.1.iv}]
After $0.044\,s$ and $7$ iterations, DO converged to the $\epsilon$-NE:
\begin{align*}
\mu_1^\star &\approx 0.333\delta_{ (1, 0, 0) } + 0.334\delta_{ (0, 1, 0) } + 0.334\delta_{ (0, 0, 1) },\\
\mu_2^\star &\approx 1\delta_{ (0.33, 0.33, 0.33) }.
\end{align*}
\item[Result of Example \ref{ex.nie21.6.2.i}]
After $0.009\,s$ and $2$ iterations, DO converged to the $\epsilon$-NE:
\begin{align*}
x_{1}^\star \approx  (0, 0) , \;
x_{2}^\star \approx  (1, 0) .
\end{align*}
\item[Result of Example \ref{ex.nie21.6.2.ii}]
After $0.047\,s$ and $8$ iterations, DO converged to the $\epsilon$-NE:
\begin{align*}
\mu_1^\star &\approx 0.289\delta_{ (0, 0, 0) } + 0.711\delta_{ (0, 0, 1) },\\
\mu_2^\star &\approx 0.376\delta_{ (0.69, 0.69, 1) } + 0.624\delta_{ (0.72, 0.72, 1) }.
\end{align*}
\item[Result of Example \ref{ex.nie21.6.3.i}]
After $0.03\,s$ and $3$ iterations, DO converged to the $\epsilon$-NE:
\begin{align*}
x_{1}^\star \approx  (-1, 1, -1) , \;
x_{2}^\star \approx  (1, 1, 1) .
\end{align*}
\item[Result of Example \ref{ex.nie21.6.3.ii}]
After $0.029\,s$ and $5$ iterations, DO converged to the $\epsilon$-NE:
\begin{align*}
\mu_1^\star &\approx 0.12\delta_{ (1, 1, -1) } + 0.16\delta_{ (1, -1, -1) } + 0.16\delta_{ (1, 1, 1) } + 0.28\delta_{ (-1, -1, -1) } + 0.28\delta_{ (-1, 1, 1) },\\
\mu_2^\star &\approx 0.12\delta_{ (1, 1, -1) } + 0.16\delta_{ (1, -1, -1) } + 0.16\delta_{ (1, 1, 1) } + 0.28\delta_{ (-1, -1, -1) } + 0.28\delta_{ (-1, 1, 1) }.
\end{align*}
\item[Result of Example \ref{ex.nie21.6.4.i}]
After $0.061\,s$ and $2$ iterations, DO converged to the $\epsilon$-NE:
\begin{align*}
\mu_1^\star &\approx \delta_{ (-1, 0, 0) },\\
\mu_2^\star &\approx 0.5\delta_{ (0, 0, 1) } + 0.5\delta_{ (1, 0, 0) }.
\end{align*}
\item[Result of Example \ref{ex.nie21.6.4.ii}]
After $1.5\,s$ and $13$ iterations, DO converged to the $\epsilon$-NE:
\begin{align*}
\mu_1^\star &\approx 0.078\delta_{ (0.71, -0.71, 0.00) } + 0.098\delta_{ (0.53, -0.8, 0.27) } + 0.103\delta_{ (-0.81, 0.46, 0.35) } \\ &+ 0.104\delta_{ (-0.77, 0.13, 0.63) } + 0.107\delta_{ (-0.58, 0.79, -0.21) } + 0.113\delta_{ (-0.32, 0.81, -0.49) } \\ &+ 0.121\delta_{ (-0.07, -0.67, 0.74) } + 0.136\delta_{ (0.69, 0.03, -0.72) } + 0.14\delta_{ (-0.45, -0.36, 0.82) },\\
\mu_2^\star &\approx 0.008\delta_{ (0.58, 0.61, 0.54) } + 0.209\delta_{ (0.61, 0.57, 0.56) } + 0.783\delta_{ (-0.57, -0.58, -0.58) }.
\end{align*}
\item[Result of Example \ref{ex.nie21.6.5}]
After $0.094\,s$ and $10$ iterations, DO converged to the $\epsilon$-NE:
\begin{align*}
\mu_1^\star &\approx 0.204\delta_{ (0.74, 0.44, 0.37) } + 0.796\delta_{ (0.72, 0.47, 0.34) },\\
\mu_2^\star &\approx 0.5\delta_{ (0.7, 0.57, 0.44) } + 0.5\delta_{ (0.69, 0.53, 0.49) }.
\end{align*}
\item[Result of Example \ref{ex.nie21.6.6}]
After $0.13\,s$ and $7$ iterations, DO converged to the $\epsilon$-NE:
\begin{align*}
\mu_1^\star &\approx 0.333\delta_{ (0, 0, 1) } + 0.333\delta_{ (0, 1, 0) } + 0.333\delta_{ (1, 0, 0) },\\
\mu_2^\star &\approx \delta_{ (0.58, 0.58, 0.58) }.
\end{align*}
\item[Result of Example \ref{ex.nie21.6.10}]
After $0.017\,s$ and $2$ iterations, DO converged to the $\epsilon$-NE:
\begin{align*}
\mu_1^\star &\approx \delta_{ (0, 1, 0, 0, 0) },\\
\mu_2^\star &\approx 0.167\delta_{ (0, 1, 0, 0, 0) } + 0.833\delta_{ (1, 0, 0, 0, 0) }.
\end{align*}
\end{description}


\subsection{Experiments on Two-Player Zero-Sum Rational Games}
Double Oracle recovered pure strategies in the games where they exist, except for \cref{ex.zhou19.5.5.ii,ex.zhou19.5.6.ii}. Tightening the $\epsilon$ in those cases did not produce simpler equilibria.

\begin{description}
\item[Result of Example \ref{ex.zhou19.5.3.i}]
After $0.0061\,s$ and $1$ iterations, DO converged to the $\epsilon$-NE:
\begin{align*}
x_{1}^\star \approx  (0, 0) ,\;
x_{2}^\star \approx  (0, 0) .
\end{align*}
\item[Result of Example \ref{ex.zhou19.5.3.ii}]
After $0.16\,s$ and $6$ iterations, DO converged to the $\epsilon$-NE:
\begin{align*}
\mu_1^\star &\approx 0.149\delta_{ (0, 0, 0) } + 0.425\delta_{ (0, 0, 1) } + 0.425\delta_{ (0, 0, 1) },\\
\mu_2^\star &\approx 0.46\delta_{ (0.7, 0.24, 1) } + 0.54\delta_{ (0.58, 0.24, 0) }.
\end{align*}
\end{description}

\subsection{Games from Classical Literature}
\begin{description}
\item[Result of Example \ref{ex.dresher61.pg110}]
After $0.018\,s$ and $4$ iterations, DO converged to the $\epsilon$-NE:
\begin{align*}
\mu_1^\star &\approx 0.5\delta_{ 0 } + 0.5\delta_{ 1 },\\
\mu_2^\star &\approx \delta_{ 0.5 }.
\end{align*}
\item[Result of Example \ref{ex.dresher61.pg111}] Converged in $0.14\,s$ and $8$ iterations to:
\begin{align*}
\mu_1^\star &\approx 0.523\delta_{ 1 } + 0.477\delta_{ 0 },\\
\mu_2^\star &\approx 0.919\delta_{ 0.5 } + 0.081\delta_{ 0.53 }.
\end{align*}
\item[Result of Example \ref{ex.gross58.poly}]
After $0.042\,s$ and $2$ iterations, DO converged to the $\epsilon$-NE:
\begin{align*}
x_{1}^\star \approx  0.5 ,\;
x_{2}^\star \approx  0.5 .
\end{align*}
\item[Result of Example \ref{ex.karlin59.vol2.sec71.ex1}]
After $0.074\,s$ and $6$ iterations, DO converged to the $\epsilon$-NE:
\begin{align*}
\mu_1^\star &\approx 1\delta_{ 0.5 },\\
\mu_2^\star &\approx 0.499\delta_{ 1 } + 0.501\delta_{ 0 }.
\end{align*}
\item[Result of Example \ref{ex.karlin59.vol2.sec71.ex2}]
After $0.1\,s$ and $4$ iterations, DO converged to the $\epsilon$-NE:
\begin{align*}
\mu_1^\star &\approx 0.379\delta_{ 0 } + 0.621\delta_{ 1 },\\
\mu_2^\star &\approx 0.313\delta_{ 1 } + 0.687\delta_{ 0.11 }.
\end{align*}
\item[Result of Example \ref{ex.karlin59.vol2.sec76.pr1}]
After $0.043\,s$ and $3$ iterations, DO converged to the $\epsilon$-NE:
\begin{align*}
\mu_1^\star &\approx 0.25\delta_{ 0 } + 0.75\delta_{ 1 },\\
\mu_2^\star &\approx 0.25\delta_{ 0 } + 0.75\delta_{ 1 }.
\end{align*}
\item[Result of Example \ref{ex.karlin59.vol2.sec76.pr2}]
After $0.12\,s$ and $6$ iterations, DO converged to the $\epsilon$-NE:
\begin{align*}
\mu_1^\star &\approx 0.054\delta_{ 0.72 } + 0.087\delta_{ 0 } + 0.087\delta_{ 0 } + 0.144\delta_{ 1 } + 0.628\delta_{ 0.41 },\\
\mu_2^\star &\approx 0.08\delta_{ 0 } + 0.231\delta_{ 0.41 } + 0.333\delta_{ 0.19 } + 0.356\delta_{ 0.8 }.
\end{align*}

\end{description}

\subsection{Experiments on Applied Games and Others}
\begin{description}
\item[Result of Example \ref{ex.chasnov20.5.2}]
After $0.037\,s$ and $3$ iterations, DO converged to the $\epsilon$-NE:
\begin{align*}
\mu_1^\star &\approx \delta_{ -1.08 },\\
\mu_2^\star &\approx 0.5\delta_{ 0.98 } + 0.5\delta_{ 0.98 }.
\end{align*}
\item[Result of Example \ref{ex.chasnov20.b.1}]
After $0.034\,s$ and $3$ iterations, DO converged to the $\epsilon$-NE:
\begin{align*}
\mu_1^\star &\approx 0.5\delta_{ 0.0 } + 0.5\delta_{ 1.0 },\\
\mu_2^\star &\approx 0.5\delta_{ 0.0 } + 0.5\delta_{ 1.0 }.
\end{align*}
\item[Result of Example \ref{ex.ratliff13.location}]
After $0.02\,s$ and $2$ iterations, DO converged to the $\epsilon$-NE:
\begin{align*}
x_{1}^\star \approx  0 ,\;
x_{2}^\star \approx  -3.14 .
\end{align*}
\item[Result of Example \ref{ex.kroupa21.5}]
After $0.32\,s$ and $9$ iterations, DO converged to the $\epsilon$-NE:
\begin{align*}
\mu_1^\star &\approx 0.218\delta_{ -0.08 } + 0.782\delta_{ -0.06 },\\
\mu_2^\star &\approx 0.161\delta_{ 0.37 } + 0.839\delta_{ 0.35 },\\
\mu_3^\star &\approx 0.278\delta_{ -1 } + 0.722\delta_{ 1 }.
\end{align*}
\item[Result of Example \ref{ex.kroupa21.6}]
After $1.6\,s$ and $11$ iterations, DO converged to the $\epsilon$-NE:
\begin{align*}
\mu_1^\star &\approx 0.001\delta_{ (0.3, 0, 0, 0.2) } + 0.218\delta_{ (0.29, 0, 0, 0.21) } + 0.245\delta_{ (0.29, 0, 0, 0.21) } \\ &+ 0.268\delta_{ (0.29, 0, 0, 0.21) } + 0.268\delta_{ (0.29, 0, 0, 0.21) },\\
\mu_2^\star &\approx 0.278\delta_{ (0.34, 0, 0.22, 0.24) } + 0.279\delta_{ (0.32, 0.25, 0, 0.23) } + 0.444\delta_{ (0.4, 0.12, 0, 0.28) },\\
\mu_3^\star &\approx 0.157\delta_{ (1, 0) } + 0.158\delta_{ (0.97, 0.03) } + 0.171\delta_{ (0.06, 0.94) } + 0.171\delta_{ (0.05, 0.95) } \\ &+ 0.171\delta_{ (0.02, 0.98) } + 0.171\delta_{ (0, 1) },\\
\mu_4^\star &\approx 0.239\delta_{ (0, 1) } + 0.761\delta_{ (1, 0) }.
\end{align*}
\item[Result of Example \ref{ex.adam21.townsend}]
After $0.62\,s$ and $11$ iterations, DO converged to the $\epsilon$-NE:
\begin{align*}
\mu_1^\star &\approx 0.008\delta_{ -2.09 } + 0.013\delta_{ 1.33 } + 0.979\delta_{ 0.03 },\\
\mu_2^\star &\approx 0.064\delta_{ -2.17 } + 0.137\delta_{ -0.9 } + 0.253\delta_{ 0.93 } + 0.546\delta_{ -2.11 }.
\end{align*}
\item[Result of Example \ref{ex.yasodharan19}]
After $0.015\,s$ and $3$ iterations, DO converged to the $\epsilon$-NE:
\begin{align*}
x_{1}^\star \approx  (1, 1) ,\;
x_{2}^\star \approx  0.8 .
\end{align*}
\item[Result of Example \ref{ex.nguyen23.iv}]
After $0.37\,s$ and $2$ iterations, DO converged to the $\epsilon$-NE:
\begin{align*}
x_{1}^\star \approx (0.33, 0.34, 0.34, 0.23, 0.23, 0.24).
\end{align*}
\end{description}



\subsection{Experiments on Optimization Test Functions}
As we searched for an $\epsilon$-NE with $\epsilon=0.001$ consistently across all games, two of the following examples had to be rescaled. Example \ref{ex.surjanovic.boha1} was scaled by $10^{-2}$, and Example \ref{ex.surjanovic.rosen} by $10^{-3}$ as its utility reaches over $10^6$ making the stopping criterion impractically strict.
\begin{description}
\item[Example \ref{ex.surjanovic.ackley}]
In $2$ iterations and $0.023\,s$, DO converged to the equilibrium:
\begin{align*}
x_{1} \approx { -32.5 },\;
x_{2} \approx { 0 }.
\end{align*}
\item[Example \ref{ex.surjanovic.boha1}]
In $2$ iterations and $0.012\,s$, DO converged to the equilibrium:
\begin{align*}
x_{1} \approx { -100.0 },\;
x_{2} \approx { 0 }.
\end{align*}
\item[Example \ref{ex.surjanovic.branin}]
In $11$ iterations and $0.18\,s$, DO converged to the equilibrium:
\begin{align*}
\mu_1^\star &\approx 0.151\delta_{ 6.23 } + 0.362\delta_{ 6.25 } + 0.487\delta_{ -5.0 },\\
\mu_2^\star &\approx 0.023\delta_{ 9.01 } + 0.464\delta_{ 8.91 } + 0.514\delta_{ 8.95 }.
\end{align*}
\item[Example \ref{ex.surjanovic.mccorm}]
In $10$ iterations and $0.15\,s$, DO converged to the equilibrium:
\begin{align*}
\mu_1^\star &\approx 0.403\delta_{ -1.5 } + 0.597\delta_{ 4 },\\
\mu_2^\star &\approx 0.128\delta_{ 0.51 } + 0.872\delta_{ 0.49 }.
\end{align*}
\item[Example \ref{ex.surjanovic.michal}]
In $2$ iterations and $0.015\,s$, DO converged to the equilibrium:
\begin{align*}
x_{1} \approx { 0 },\;
x_{2} \approx { 1.57 }.
\end{align*}
\item[Example \ref{ex.surjanovic.rosen}]
In $2$ iterations and $0.013\,s$, DO converged to the equilibrium:
\begin{align*}
x_{1} \approx { 10 },\;
x_{2} \approx { 10 }.
\end{align*}
\end{description}

\subsection{Results on Discontinuous Games}
Our DO algorithm is not guaranteed to converge on discontinuous games; However, as long as the discontinuities can be avoided, the algorithm often recovers a solution.
The following games by \textcite{zhou.wang.ea:saddle} contain discontinuities, but converge without any reformulation or modification to the algorithm.

\begin{description}
\item[Example \ref{ex.zhou19.5.1}]
Converged in $0.1\;(+0.002)\,s$ and $4$ iterations to:
\begin{align*}
x_1^\star \approx (0.32, 0.32, 0.37),\;
x_2^\star \approx (0, 0, 1).
\end{align*}
\item[Example \ref{ex.zhou19.5.5.i}]
Converged in $0.22\;(+0.002)\,s$ and $6$ iterations to:
\begin{align*}
\mu_1^\star &\approx 0.424\delta_{ (0.71, 0.71) } + 0.307\delta_{ (-0.71, -0.71) } + 0.269\delta_{ (-0.71, -0.71) },\\
\mu_2^\star &\approx 1.0\delta_{ (0, 0) }.
\end{align*}
\item[Example \ref{ex.zhou19.5.5.ii}]
Converged in $0.38\;(+0.002)\,s$ and $10$ iterations to:
\begin{align*}
\mu_1^\star &\approx 0.548\delta_{ (0.46, 0.47, 0.34) } + 0.452\delta_{ (0.49, 0.43, 0.35) },\\
\mu_2^\star &\approx 0.641\delta_{ (0.67, 0.54, 0.5) } + 0.359\delta_{ (0.64, 0.59, 0.49) }.
\end{align*}
\item[Example \ref{ex.zhou19.5.6.ii}]
Converged in $0.062\;(+0.006)\,s$ and $3$ iterations to:
\begin{align*}
x_1^\star \approx (0.41, 0.41, 0.81),\;
x_2^\star \approx (-0.39, -0.84, -0.39).
\end{align*}
\item[Example \ref{ex.zhou19.5.7.i}]
Converged in $0.035\;(+0.004)\,s$ and $2$ iterations to:
\begin{align*}
x_1^\star \approx (0.42, 0.64, 0.64),\;
x_2^\star \approx (0.54, 0.84).
\end{align*}
\end{description}

DO initially failed to solve Examples \ref{ex.zhou19.5.2.ii} and \ref{ex.zhou19.5.7.i.polar} as it attempted to initialize at the discontinuity. Both examples were solved by initializing in the interior of the sets.
\begin{description}
\item[Example \ref{ex.zhou19.5.2.ii}]
Converged in $0.0086\;(+0.002)\,s$ and $1$ iterations to:
\begin{align*}
x_1^\star \approx (0, 0, 0),\;
x_2^\star \approx (0, 0, 0).
\end{align*}
\item[Example \ref{ex.zhou19.5.7.i.polar}]
Converged in $0.056\;(+0.002)\,s$ and $2$ iterations to:
\begin{align*}
x_1^\star \approx (0.87, 1),\;
x_2^\star \approx 1.
\end{align*}
\end{description}

Our implementation can not be directly applied to search for solutions in Games with unbounded strategy spaces. We know of no global solver powerful enough to automatically infer bounds on the location of best responses in general games.


\end{document}

